\begin{abstract}
Generating verifiable, program-based CAD models requires exact geometric precision.
While Vision-Language Models (VLMs) demonstrate remarkable capabilities in understanding
3D topologies and generating multimodal code, they consistently fail at \emph{spatial
equivariance}---struggling with deterministic mapping rules like coordinate-plane
transformations and arc-direction inversions.
Existing generative pipelines typically rely on probabilistic Re-Act loops, forcing VLMs
to ``trial-and-error'' analytical geometry parameters.
This severely increases workflow complexity without guaranteeing convergence, even when the
VLM is provided with multi-modal visual feedback.

In this paper, we propose CADLoop, an elegant neuro-symbolic data curation
pipeline that acts as a deterministic fallback mechanism outside the standard Re-Act loop.
Decoupling structural reasoning from explicit constraint execution, CADLoop equips the VLM
agent with a unified equivariant-aware skill.
The VLM identifies topological intent, while this single programmatic tool analytically
verifies and resolves complex continuous constraints (\eg, volumetric anomalies,
inner-hole overcuts) via Abstract Syntax Tree (AST) manipulation.
Our core insight is the identification of a proximal boundary at
$\text{IoU} \approx 0.8$: above this threshold, topological structures are sound, and
failures are strict single-parameter mismatches perfectly suited for automated tool-use
repair.
Targeting this $>0.8$ ``near-miss'' regime provides a high-density, low-noise signal
optimal for Supervised Fine-Tuning (SFT).
On a rigorous evaluation subset of 200 near-miss cases, baseline Re-Act resolves only
${\sim}40\%$, whereas CADLoop's neuro-symbolic fallback achieves a 100\% repair
rate.
To accelerate future multimodal 3D reasoning, we release a large-scale verified dataset
where over 95\% of pairs achieve a perfect $\text{IoU} = 1.0$, reducing the mean
Chamfer Distance by over 90\%.
\end{abstract}
